\subsection{Pass 10961: exact Albert $Cl(9)$, doubled $Cl(10)$, and the clock grade obstruction}

The nine spatial Peirce-zero Jordan multiplication operators act on the
Albert Peirce $16$ as exact rational Euclidean Clifford generators:
\[
\{\gamma_i,\gamma_j\}=2\delta_{ij}I_{16}.
\]
Their bivectors span dimension $36$.  On two copies of the Albert $16$,
the block matrices
\[
\Gamma_i=\begin{pmatrix}0&\gamma_i\\\gamma_i&0\end{pmatrix},
\qquad
\Gamma_{10}=\operatorname{diag}(I_{16},-I_{16})
\]
satisfy $\{\Gamma_a,\Gamma_b\}=2\delta_{ab}I_{32}$ and have a
$45$-dimensional bivector span, giving an explicit $Cl(10)$ Dirac system.

For the minimal doubled clock
$G=\operatorname{diag}(\zeta_8^{-1}U,\zeta_8U)$,
\[
G^2=i\Gamma_{10},\qquad G^4=-I_{32},
\]
but one tick obeys
$G\Gamma_iG^{-1}=-i\Gamma_{10}\sum_j(C_9)_{ji}\Gamma_j$ while fixing
$\Gamma_{10}$.  Thus nine Clifford vectors are sent to mixed bivectors;
the vector space and its one-tick image span dimension $19$ and intersect
only in $\Gamma_{10}$.  The full order-eight clock is therefore not a
Spin(10)/Pin(10) vector normalizer in this canonical Albert extension.
Two ticks do return to the vector normalizer, acting as
$\operatorname{diag}(-I_9,+1)$ on the ten Clifford vectors.
